% !PRINTED 173
\weberpage{173}

\[
\mu_1 + \mu_2 + \cdots + \mu_m = \nu_1 + 2\nu_2 + \cdots + (n-1)\nu_{n-1},
\]
and accordingly the expression for $\sigma_m$ can be represented thus:
\weq{15}{
\sigma_m = \overset{\mu}{\textstyle\sum} s_{\mu_1+\mu_2+\ldots+\mu_m}\, \alpha_{\mu_1}\alpha_{\mu_2}\cdots\alpha_{\mu_m}.
}
For example, we obtain
\[
\sigma_1 = \alpha_0 s_0 + \alpha_1 s_1 + \cdots + \alpha_{n-1} s_{n-1}
\]
\weq{16}{
\begin{aligned}
\sigma_2 &= \overset{i,k}{\textstyle\sum} s_{i+k}\, \alpha_i \alpha_k \\
\sigma_3 &= \overset{h,i,k}{\textstyle\sum} s_{h+i+k}\, \alpha_h \alpha_i \alpha_k,
\end{aligned}
}
in which $h$, $i$, $k$ run through the sequence of numbers $0, 1, 2 \ldots n-1$.

\begin{center}
\S.\,53.

\gesperrt{Application to the cubic and biquadratic equations.}
\end{center}

The goal that \gesperrt{Tschirnhausen} already had in view with this transformation consisted in making so many terms of the transformed equation $\Phi(y) = 0$ vanish, by determining the substitution coefficients $\alpha_0, \alpha_1 \ldots \alpha_{n-1}$, that the equation becomes solvable. This succeeds easily for the equations of the third and fourth degree, though it is not simple, and does not seem possible without lengthy calculations, to derive the otherwise well-known formulas in this way.

Let us first take the cubic equation
\weq{1}{
x^3 + a x^2 + b x + c = 0,
}
which we transform by means of the following substitution:
\weq{2}{
y = \alpha + \beta x + \gamma x^2
}
In order to reduce the equation for $y$
\weq{3}{
y^3 + A y^2 + B y + C = 0
}
to a pure cubic equation, we have to determine the ratios $\alpha:\beta:\gamma$ from the two equations
\weq{4}{
A = 0, \quad B = 0
}
of which the first is linear, the second quadratic.

The equations (4) are equivalent to
\weq{5}{
\sigma_1 = 0, \quad \sigma_2 = 0
}
%CONT>>
